% OEM engineering brief -- what survived testing.
% Build: latexmk -pdf -outdir=build oem-brief.tex
% Every number traces to a versioned artifact under data/processed/ and to a
% pre-registration tag; nothing here is hand-copied (norm N4).
\documentclass[11pt]{article}
\usepackage[a4paper,margin=2.0cm]{geometry}
\usepackage[T1]{fontenc}
\usepackage[utf8]{inputenc}
\usepackage{lmodern}
\usepackage{microtype}
\usepackage{booktabs}
\usepackage{amsmath}
\usepackage{xcolor}
\usepackage[hidelinks]{hyperref}
\usepackage{titlesec}
\usepackage{enumitem}
\setlist[itemize]{leftmargin=1.1em,itemsep=1pt,topsep=2pt}
\titlespacing*{\section}{0pt}{9pt plus 2pt}{3pt}
\titleformat{\section}{\large\bfseries}{}{0pt}{}
\titlespacing*{\subsection}{0pt}{6pt plus 1pt}{2pt}
\titleformat{\subsection}{\normalsize\bfseries}{}{0pt}{}
\setlength{\parskip}{3pt}
\setlength{\parindent}{0pt}
\definecolor{inkc}{HTML}{212529}
\definecolor{bluec}{HTML}{4263EB}
\definecolor{redc}{HTML}{A61E4D}
\newcommand{\no}{\textcolor{redc}{\textbf{No}}}
\newcommand{\yes}{\textcolor{bluec}{\textbf{Yes}}}
\pagestyle{plain}

\begin{document}

\begin{center}
{\LARGE\bfseries Which AI claims in engine performance survive testing?\par}
\vspace{1mm}
{\large A measurement rig, and what it found --- including about our own numbers\par}
\vspace{2mm}
{\small\itshape Engineering brief. All results from a pre-registered testbed on a
CFM56-7B26 model built entirely from public certification data (TCDS, ICAO EEDB).
No proprietary inputs.\par}
\end{center}
\vspace{1mm}

\hrule
\vspace{2mm}
\textbf{Summary.} We built a testbed where the true engine condition is known, so
a diagnosis can be scored against the answer rather than against another
estimate. We then ran the comparison everyone quotes and nobody audits: tuned AI
against a tuned classical gas-path stack, on identical data, with hypotheses
frozen before the runs. \textbf{Most of the advertised advantage did not
survive}, including two of our own headline numbers. What did survive is small,
specific, and --- for an OEM --- points at instrumentation rather than software.
\vspace{2mm}
\hrule

\section{1. What failed, and why it matters}

\subsection{A published deep-learning result, replicated}
A 2026 Springer paper reports a bidirectional LSTM as the best of fifteen
architectures for remaining-life prediction on NASA's C-MAPSS benchmark. We
re-implemented it exactly.

\begin{itemize}
  \item It \emph{does} replicate: $14.56$ cycles RMSE against a reported
        $14.12$. This is not a straw man.
  \item Its architecture ranking rests on gaps of $0.12$ to $0.60$ cycles.
        \textbf{The run-to-run spread from the random seed alone is $2.94$
        cycles.} The ranking is noise. The paper reports one run per model.
  \item The published architecture \textbf{fails to train on 5 of 10 runs},
        collapsing to predicting zero. The cause is specified in the paper's own
        tables: a ReLU activation on the output layer, plus a scheduler that cuts
        the learning rate hundredfold to a floor it cannot recover from.
\end{itemize}

\textbf{Takeaway for evaluation:} ask any vendor or internal team for the seed
spread before you look at the headline. A ranking inside the noise is not a
result, and a model that silently fails half the time is not a product.

\subsection{Two of our own headline numbers, corrected downward}
We audited our own claims the same way and both shrank.

\begin{itemize}
  \item \emph{``AI gives $6\times$ tighter uncertainty intervals.''} False as
        stated. Our AI interval was statistically calibrated and the classical
        one was not. Calibrate both the same way --- a technique the classical
        side can adopt for free --- and the advantage is $\mathbf{1.34\times}$.
        About three quarters of that headline was the calibration method, not
        the model.
  \item \emph{``AI predicts remaining life $2.3$--$4.4\times$ better.''} True
        only against the linear extrapolation operators field today. Against a
        competent classical prognostic it is $\mathbf{1.34\times}$ at $90\,\%$ of
        life and a \textbf{statistical tie} at $70\,\%$. The advantage is real,
        late-life, and modest.
\end{itemize}

\subsection{Where AI lost outright}
Fault isolation between confusable gas-path faults: \textbf{31\,\% versus
31\,\%}. Identical. The limit is informational, not algorithmic --- three cockpit
measurements against ten health parameters, with two fault signatures
$1.3^{\circ}$ apart. We confirmed it four independent ways: halving every sensor
noise leaves it at $0.31$--$0.38$; a ``virtual sensor'' adds \emph{exactly}
nothing; sweeping the entire calibrated operating envelope leaves the angle at
$1.38^{\circ}$; and no learned method moved it. \textbf{No algorithm removes a
rank deficiency.}

Also refuted: ``physics-informed'' learning, twice, when applied to remaining
life. And $87\,\%$ of mid-life remaining-life uncertainty is \emph{irreducible} ---
two engines in identical measured condition genuinely fail hundreds of cycles
apart. Most prognostic R\&D spend is chasing that, not a method gap.

\section{2. What survived}

\subsection{An instrumentation design tool --- and it is not AI}
The strongest result is classical estimation theory applied through the engine
deck: accumulate Fisher information over an engine's actual flown conditions and
you get, per engine and per health direction, a bound on what any diagnosis can
possibly know. Validated against the true component state, which only a
synthetic-truth testbed permits.

It answers the question an OEM actually owns --- and one caveat travels with it, because we
found it ourselves and late: the \emph{ranking} claim (that the bound orders the estimator's
true error, $\rho = 0.70$) does not clear a physics-free null. About a quarter of it is the
matrix the bound shares with the estimator; about half survives against an estimator that never
sees that matrix; and neither reaches significance, because a rank test over ten directions
cannot. The bound is not shown to be wrong --- it is shown to be unprovable the way we tested
it --- and we know precisely what
would test it, which is a fleet whose engines fly genuinely different missions. Ours vary
by $1.15\,\%$ in the quantity that had to vary, so the sharper statement is not
\emph{unproven} but \emph{unmeasurable on this fleet}. The \emph{acquisition} figures below
are separate measurements and are unaffected.

\begin{center}\small
\begin{tabular}{@{}lr@{}}
  \toprule
  Instrumentation decision & Measured effect \\
  \midrule
  Add station probes (P25/T25/PS3/T3) & HPC efficiency recoverable $45\times$ better \\
  Same probes, confusable isolation & $0.15 \rightarrow 0.92$ \\
  Halve noise on existing cockpit probes & no effect on the wall \\
  Add a ``virtual sensor'' / better algorithm & no effect (rank unchanged) \\
  \bottomrule
\end{tabular}
\end{center}

\textbf{Product-strategy consequence.} Any EHM service built on cockpit-only data
sits under a hard informational ceiling that no software roadmap removes. The way
through is instrumentation --- which an OEM controls and a customer does not.
This tool prices that decision before the hardware is committed, and it runs on
our own decks, where it is stronger than on the public model used here.

\subsection{The one AI result that held up}
Injecting physics into a learner \emph{hurts} remaining-life prediction (refuted
twice) but \textbf{helps module attribution} --- which mechanism is consuming the
engine --- by a large, significant margin. We ran it with a control arm so the
gain could not be attributed to simply adding more inputs.

The reconciliation is mechanical: remaining life is already carried by the
EGT-margin signal the model can see, so a physical state estimate adds nothing;
\emph{which module} is exactly what that estimate encodes. \textbf{``Physics-informed''
is neither a free win nor a dead end --- the task decides}, and this tells us
which of our own EHM products should carry a physics channel and which should not.

\subsection{A related finding on the wall}
The isolation wall is a property of a \emph{single snapshot}. Over a full
history, three of five degradation mechanisms become attributable from the shape
of the trajectory --- including the hot-section mechanism that drives the
confusable pair. The governing rule we measured: \textbf{trajectory shape
separates what has shape}. A wash sawtooth and a break-in knee are readable; a
featureless linear ramp is not, however long you watch.

\section{3. What this is, and what it is not}

\textbf{Not:} a product, a validated tool for operational use, or evidence about
absolute numbers. The engine model is calibrated to public certification data,
not to our test-cell data, and the fleet is synthetic. Statistical predictions
here are not certified conclusions.

\textbf{Is:} a reusable method for deciding which performance-AI claims to fund
and which to reject, with the discipline that makes the answers trustworthy ---
hypotheses frozen before runs, equal tuning budgets for both sides, negative
results reported with the same weight as positive ones, and every figure
regenerated from data rather than transcribed.

\textbf{Proposed next step.} Re-run the instrumentation analysis on our own
engine decks and our own candidate sensor suites. It requires no AI, no new data
collection, and answers a question we already pay to answer by other means.

\vfill
{\footnotesize\itshape Every number above traces to a versioned artifact and a
pre-registration tag in the study repository; the supporting document runs to 135
pages with the full protocol, the refutations, and the audit trail.\par}

\end{document}
